% from ega2-1-fr.tex
\oldpage[II]{163}
tels que $\varepsilon$ et $j$ soient des immersions fermées, et $i$ un
morphisme affine, qui est une immersion ouverte dominante, pour laquelle
\begin{equation}
\label{II.8.3.2.3-fr}
  i(C)=\widehat{C}-j(X)~;
\tag{8.3.2.3}
\end{equation}
en outre $\widehat{C}$ est le plus petit sous-préschéma fermé de $\widehat{C}$
majorant $i(C)$.
\end{proposition}

Pour définir $i$, on considère l'ouvert de $\widehat{C}$,
\begin{equation}
\label{II.8.3.2.4-fr}
  \widehat{C}_z=
  \operatorname{Spec}(\widehat{\mathcal{S}}/(z-1)\widehat{\mathcal{S}})
\tag{8.3.2.4}
\end{equation}
(\hyperref[II.3.1.4-fr]{3.1.4}), $z$ s'identifiant canoniquement à une
section de $\widehat{\mathcal{S}}$ au-dessus de $Y$. L'isomorphisme
$i:C\xrightarrow{\sim}\widehat{C}_z$ correspond alors à l'isomorphisme
canonique (\hyperref[II.8.2.3.1-fr]{8.2.3.1})
\[
  \widehat{\mathcal{S}}/(z-1)\widehat{\mathcal{S}}
  \xrightarrow{\sim}\mathcal{S}.
\]
Le morphisme $\varepsilon$ correspond à l'homomorphisme d'augmentation
$\mathcal{S}\to\mathcal{S}_0=\mathcal{O}_Y$ ayant pour noyau
$\mathcal{S}_+$ (\hyperref[II.1.2.7-fr]{1.2.7}) et comme ce dernier est
surjectif, $\varepsilon$ est une immersion fermée
(\hyperref[II.1.4.10-fr]{1.4.10}). Enfin, $j$ correspond de même
(\hyperref[II.3.5.1-fr]{3.5.1}) à l'homomorphisme surjectif de degré $0$,
$\widehat{\mathcal{S}}\to\mathcal{S}$, qui se réduit à l'identité dans
$\mathcal{S}$ et est nul dans $z\widehat{\mathcal{S}}$, qui est son noyau~;
$j$ est partout défini et est une immersion fermée en vertu de
(\hyperref[II.3.6.2-fr]{3.6.2}).

Pour démontrer les autres assertions de
(\hyperref[II.8.3.2-fr]{8.3.2}), on peut évidemment se borner au cas où
$Y=\operatorname{Spec}(A)$ est affine, $\mathcal{S}=\widetilde{S}$, où $S$
est une $A$-algèbre graduée, d'où
$\widehat{\mathcal{S}}=(\widehat{S})^{\sim}$~; les éléments homogènes $f$ de
$S_+$ s'identifient alors à des sections de $\widehat{\mathcal{S}}$ au-dessus
de $Y$, et l'ouvert de $\widehat{C}$ noté $D_+(f)$ dans
(\hyperref[II.2.3.3-fr]{2.3.3}) s'écrit aussi $\widehat{C}_f$
(\hyperref[II.3.1.4-fr]{3.1.4})~; de même l'ouvert de $C$ noté $D(f)$ dans
(\hyperref[I.1.1.1-fr]{I, 1.1.1}) s'écrit aussi $C_f$
(\hyperref[0.5.5.2-fr]{0, 5.5.2}). Cela étant, il résulte de
(\hyperref[II.2.3.14-fr]{2.3.14}) et de la définition de $\widehat{S}$ que
dans le cas présent, les ouverts $\widehat{C}_z=i(C)$ et $\widehat{C}_f$
($f$ homogène dans $S_+$) constituent un \emph{recouvrement} de
$\widehat{C}$. En outre, on a, avec ces notations,
\begin{equation}
\label{II.8.3.2.5-fr}
  i^{-1}(\widehat{C}_f)=C_f~;
\tag{8.3.2.5}
\end{equation}
en effet,
$\widehat{C}_f\cap i(C)=\widehat{C}_f\cap\widehat{C}_z
=\widehat{C}_{fz}=\operatorname{Spec}(\widehat{S}_{(fz)})$. Or, si
$d=\deg(f)$, $\widehat{S}_{(fz)}$ est canoniquement isomorphe à
$(\widehat{S}_{(z)})_{f/z^d}$ (\hyperref[II.2.2.2-fr]{2.2.2}), et il résulte
de la définition de l'isomorphisme
(\hyperref[II.8.2.3.1-fr]{8.2.3.1}) que l'image de
$(\widehat{S}_{(z)})_{f/z^d}$ par l'isomorphisme correspondant des anneaux de
fractions est exactement $S_f$. Comme $C_f=\operatorname{Spec}(S_f)$, cela
prouve (\hyperref[II.8.3.2.5-fr]{8.3.2.5}) et montre en même temps que le
morphisme $i$ est affine~; en outre, la restriction de $i$ à $C_f$,
considérée comme morphisme dans $\widehat{C}_f$, correspond
(\hyperref[I.1.7.3-fr]{I, 1.7.3}) à l'homomorphisme canonique
$\widehat{S}_{(f)}\to\widehat{S}_{(fz)}$ et en vertu de ce qui précède et de
(\hyperref[II.8.2.3.2-fr]{8.2.3.2}), on peut énoncer le résultat suivant~:

\begin{env}[8.3.2.6]
\label{II.8.3.2.6-fr}
Si $Y=\operatorname{Spec}(A)$ est affine, et $\mathcal{S}=\widetilde{S}$,
alors, pour tout $f$ homogène dans $S_+$, $\widehat{C}_f$ s'identifie
canoniquement à $\operatorname{Spec}(S_f^{\leq})$, et le morphisme
$C_f\to\widehat{C}_f$ restriction de $i$ correspond alors à l'injection
canonique $S_f^{\leq}\to S_f$.
\end{env}

Notons maintenant que (pour $Y$ affine) le complémentaire de
$\widehat{C}_z$ dans $\widehat{C}=\operatorname{Proj}(\widehat{S})$

